\chapter{Transcripciones de las entrevistas}
\label{anexo:entrevistas}

Criterio de depuración. Las siguientes transcripciones se presentan en formato de diálogo directo, en primera persona para entrevistadores y profesionales entrevistados. Se eliminaron pruebas técnicas de cámara o audio, repeticiones sin valor analítico, muletillas excesivas y errores evidentes de la transcripción automática. No se reemplazaron las respuestas por síntesis en tercera persona, ya que el anexo funciona como registro primario del relevamiento cualitativo. Cuando una frase del audio resultaba dudosa, se conservó el sentido más probable sin agregar conclusiones externas.

Identificación de participantes. Los nombres e instituciones se consignan según la información provista por los autores del PFI para esta versión de trabajo. El cuerpo de cada entrevista mantiene el intercambio en primera persona y no incluye análisis interpretativo; dicho análisis se desarrollará en el capítulo de user research. % CROSS_REFERENCE_MAPPING_REQUIRED

% SOURCE_NUMBERING_ANOMALY_PRESERVED: E.1 literal in DOCX Heading2
\phantomsection
\section*{E.1 Entrevista 1 (EN-01) - Dra. Susana Torres (M.N. 179537)}
\addcontentsline{toc}{section}{E.1 Entrevista 1 (EN-01) - Dra. Susana Torres (M.N. 179537)}
\label{anexo:entrevista-en01}

\noindent\textbf{Dra. Susana Torres:} Ustedes serían la parte de programación específica, o sea, desarrollar la idea. ¿Qué es lo que necesitan? ¿Ciertos parámetros? ¿Qué tienen que tener en cuenta al momento de evaluar una persona en columna lumbar? Más que nada, si es para estenosis, cuántas patologías abarcarían. Me pregunté un montón de cosas.

\noindent\textbf{Enzo Asplanatti:} Por ahora nosotros mantenemos la parte ingenieril, que sería automatizar ese proceso y, con ayuda de la computadora, lograr encontrar la mayor cantidad de hallazgos posibles. Todavía no tenemos el alcance completamente acotado a patologías específicas. Por ahora lo estamos achicando a esto.

\noindent\textbf{Francisco Fabrello:} Estamos esperando la respuesta de la cátedra. Presentamos el proyecto, nos dijeron que avancemos y que después nos iban a dar correcciones finas. Entonces, por ahora, la idea es entrenar el modelo con datos ya etiquetados para que aprenda a detectar, y después llevarlo hacia lo que nos recomienden. También nos pidieron validar con profesionales, haciendo entrevistas con algunas preguntas, para ver qué opinan y qué podríamos agregar o quitar.

\noindent\textbf{Dra. Susana Torres:} ¿Ustedes tienen una lista de preguntas?

\noindent\textbf{Francisco Fabrello:} Sí, más o menos. La primera sería: ¿cómo trabajás con resonancias y qué lugar ocupan dentro de tu práctica?

\noindent\textbf{Dra. Susana Torres:} Resonancias lumbares, hoy por hoy estoy trabajando en FLENI. El principal volumen son los cerebros, pero después le sigue seguro la resonancia lumbar por hernias discales. Hay otros hallazgos no pensados dentro del diagnóstico presuntivo, pero la mayoría es por radiculopatía o compresión radicular.

\noindent\textbf{Francisco Fabrello:} Entonces, ¿cuál sería más o menos el caso que encontrás más frecuente?

\noindent\textbf{Dra. Susana Torres:} Protrusiones discales.

\noindent\textbf{Francisco Fabrello:} ¿Y eso sería, en palabras más simples?

\noindent\textbf{Dra. Susana Torres:} La hernia discal es como la entidad general. Hay distintos grados: está el abombamiento, la protrusión y la extrusión. La protrusión es la más frecuente.

\noindent\textbf{Francisco Fabrello:} Perfecto. Y cuando tenés que hacer el análisis de una resonancia, ¿cómo sería el recorrido para visualizarla?

\noindent\textbf{Dra. Susana Torres:} Lo que ya viene estudiado, que también hace poco se usa en FLENI, es más que nada la parte de volumetría encefálica, que tira resultados de cuantificación de sustancia gris, sustancia blanca, etcétera; pero para resonancia lumbar todavía no he trabajado con eso. Mi orden de trabajo generalmente es dividir la pantalla. El axial T2 es el que siempre estoy corroborando y en la otra pantalla veo el sagital T2, el sagital T1 y el sagital STIR para ver si hay edema óseo o cambios degenerativos a nivel de los cuerpos vertebrales. Una vez visto eso, los espacios y la alineación, paso al sagital T2 y al axial T2. En el axial puedo ver si hay alguna compresión, cómo está el canal, si hay alguna estrechez a nivel del canal y cómo están los neuroforámenes.

\noindent\textbf{Francisco Fabrello:} Perfecto. ¿Qué considerás más importante para que el informe sea útil?

\noindent\textbf{Dra. Susana Torres:} Tener en cuenta la clínica del paciente y el diagnóstico presuntivo. Eso es lo que le está doliendo. Puede haber un montón de cambios degenerativos, pero si el paciente viene por una hernia L5-S1, voy enfocada ahí. Ahí hago mi conclusión: dónde le duele y por qué le duele.

\noindent\textbf{Francisco Fabrello:} Si este sistema pudiera delimitar automáticamente las vértebras, los discos, el canal espinal y, a partir de eso, calcular mediciones como la altura de discos, el diámetro AP del canal o la alineación, ¿qué utilidad real podría tener o qué dudas te generaría?

\noindent\textbf{Dra. Susana Torres:} A eso que nombraste le agregaría, además de la altura de los discos, la intensidad de señal de los discos. Pueden tener altura respetada, pero ya verse signos de deshidratación, que es súper común y así empiezan los cambios degenerativos. El diámetro AP del canal, en cuanto a lumbar, sería más que nada el área. Muchas veces hay cambios degenerativos en los elementos posteriores y facetas que se hipertrofian. Eso genera una disminución del diámetro transverso del canal. Entonces, por ahí el AP está bien, pero como tenés cambios degenerativos atrás, las raíces de la cola de caballo se apiñan y eso ya es una estenosis leve. También hay que tener en cuenta que el paciente está acostado; con ciertos movimientos o parado puede cambiar. En cuanto a la alineación, me parece re importante, porque el paciente está acostado y no tiene la sobrecarga de la gravedad. La alineación posterior también es súper importante para sospechar algún tipo de lisis o listesis.

\noindent\textbf{Francisco Fabrello:} Si el sistema trabajara principalmente con cortes sagitales T1 y T2, pero no con cortes axiales, ¿qué podría evaluarse y qué quedaría incompleto o riesgoso?

\noindent\textbf{Dra. Susana Torres:} Bastante quedaría incompleto. Al hacer el análisis nivel por nivel, si me hacés elegir, es el sagital T2 y el axial T2. Ahí veo también la extensión hacia los neuroforámenes. Incluir la estenosis a nivel de los neuroforámenes es re importante y se ve mejor correlacionando los dos planos.

\noindent\textbf{Francisco Fabrello:} ¿Qué errores de una herramienta de este tipo serían graves o inaceptables desde el punto de vista clínico?

\noindent\textbf{Dra. Susana Torres:} Los mismos que tenemos nosotros, básicamente: sobredimensionar algo o infravalorar algún hallazgo. Lo importante está todo en la calidad de los datos que le carguen, básicamente. Y qué ground truth es la que utilizan.

\noindent\textbf{Francisco Fabrello:} Claro. Yendo más al lado de la interfaz, ¿qué tendría que mostrar o permitir el sistema para que resulte confiable y no agregue ruido o haga más tediosa la rutina?

\noindent\textbf{Dra. Susana Torres:} Pienso en lo que evalúo a diario. Por ejemplo, si me dice que tiene una estenosis grado 2 de los neuroforámenes, poder verlo. Que me lo marque primero y, si estoy de acuerdo, lo dejo; si no, hago una corrección manual y digo que es grado 1 o grado 2. Igual ahí también hay sesgo, porque si una IA te dice que es grado 2 y quizás es grado 1, depende de la experiencia y de cómo se interprete el software.

\noindent\textbf{Francisco Fabrello:} Para que este software sea viable desde el punto de vista clínico, ¿contra qué tendríamos que validarlo? Por ejemplo, segmentaciones manuales de especialistas o mediciones de expertos.

\noindent\textbf{Dra. Susana Torres:} Sí, segmentación manual de especialistas sería ideal. Yo trabajo para una empresa que se llama Prenuvo, que hace total body segmentations con resonancias y distintos tipos de informes para hallazgos positivos. Las segmentaciones las hacemos entre médicos, pero siempre hay como un pool; cada estudio está tres veces chequeado. La calidad de los datos me parece lo más importante. Se nota cuando el que primero hizo la revisión es un médico general, un residente de radiología o un especialista.

\noindent\textbf{Francisco Fabrello:} ¿Qué alcance te parecería más defendible? ¿Qué sacarías para no prometer de más?

\noindent\textbf{Dra. Susana Torres:} Lo acotaría. Como les preguntaba al principio: ¿cuál era el objetivo? ¿Valorar estenosis? ¿Delimitar estenosis de canal y estenosis neuroforaminal? Porque si no hay un montón de patologías y un montón de cuestiones que se pueden evaluar. Generalmente, si me preguntás qué es lo que más vemos y lo que generalmente causa dolor o consulta del paciente, son cambios degenerativos. También cambios tipo Modic de las plataformas. Lo acotaría, porque si no va a ser muy difícil abarcar todo y se les va a hacer eterno.

\noindent\textbf{Francisco Fabrello:} ¿Alguna red flag que te haría decir que el proyecto no debería plantearse así?

\noindent\textbf{Dra. Susana Torres:} No veo que tuviera alguna red flag. La verdad es un tema bastante estudiado. Ya hay distintos modelos e instituciones afuera que están trabajando así, con informes y resultados estructurados. Ayudaría un montón. Recalco nuevamente la calidad de los datos de entrenamiento, pero el resto me parece fantástico.

\noindent\textbf{Francisco Fabrello:} ¿Qué dirías que sería un diferencial del proyecto? Por ejemplo, que genere el informe y muestre paso a paso lo que hizo, o que deje visualizar la resonancia con el análisis encima.

\noindent\textbf{Dra. Susana Torres:} Visualmente, creo que hablo por la mayoría de los radiólogos: nosotros nos guiamos mucho por lo visual. Si las segmentaciones están como máscaras dentro de los hallazgos y yo veo que tiene pérdida de altura del disco y me muestran dónde lo está viendo, eso me genera confianza. Lo mismo con las estenosis, ya sea del canal o de los forámenes.

\noindent\textbf{Francisco Fabrello:} ¿Qué opinás de que el sistema permita cargar una resonancia anterior y una actual para hacer una comparativa?

\noindent\textbf{Dra. Susana Torres:} Buenísimo. Eso ayudaría a decir qué era lo que le dolía antes, qué aumentó y cómo fueron evolucionando esos cambios degenerativos. Sumaría un montón.

\noindent\textbf{Enzo Asplanatti:} ¿Cómo continúa después, si avanza el proyecto?

\noindent\textbf{Dra. Susana Torres:} Si el proyecto continúa y progresa, yo me dedico a segmentación y control de calidad, así que tienen mi contacto. Me interesa mucho la inteligencia artificial porque es lo que se viene. Cualquier cosa que necesiten, tienen mi contacto.

\noindent\textbf{Francisco Fabrello:} La idea es que primero validemos con profesionales la viabilidad y después, más adelante, con un producto semiarmado, podamos volver a hacer preguntas tipo checklist del uno al cinco para evaluar qué mejorar.

\noindent\textbf{Dra. Susana Torres:} Claro, cuál es la practicidad en el día a día.

\noindent\textbf{Francisco Fabrello:} También pensamos que el sistema muestre la máscara encima, una tabla de datos y que el profesional pueda corregir el dato que predijo, porque vimos que el modelo a veces falla.

\noindent\textbf{Dra. Susana Torres:} O sea, que pasa.

\noindent\textbf{Francisco Fabrello:} Por ejemplo, mide el canal, lo pinta, pero le falta un pedacito. Entiendo que depende de la resonancia, de si es T1 o T2, y de qué estructuras se ven mejor.

\noindent\textbf{Dra. Susana Torres:} Claro. En cuanto a lo lumbar, les recalco que no solamente importa el diámetro, sino también la pérdida de líquido entre las raíces. A nivel cervical puede ser solamente diámetro, pero a nivel lumbar tenés la cola de caballo. Entonces hay que medir cuán juntas están las raíces y si se sigue viendo líquido entre ellas.

\noindent\textbf{Francisco Fabrello:} Lo vamos a tener en cuenta para entrenar la IA.

\noindent\textbf{Enzo Asplanatti:} Eso está bueno.

\noindent\textbf{Dra. Susana Torres:} Bueno, muchas gracias.

\noindent\textbf{Francisco Fabrello:} Gracias. ¿En un futuro te podríamos escribir para seguir validando?

\noindent\textbf{Dra. Susana Torres:} Sí, encantada, ningún problema. Me interesa mucho.

\noindent\textbf{Francisco Fabrello:} Muchas gracias.

\noindent\textbf{Dra. Susana Torres:} Dale, nos vemos.

\noindent\textbf{Enzo Asplanatti:} Chau, chau.

% SOURCE_NUMBERING_ANOMALY_PRESERVED: E.2 literal in DOCX Heading2
\phantomsection
\section*{E.2 Entrevista 2 (EN-02) - Dr.~Gonzalo Brito}
\addcontentsline{toc}{section}{E.2 Entrevista 2 (EN-02) - Dr.~Gonzalo Brito}
\label{anexo:entrevista-en02}

\noindent\textbf{Enzo Asplanatti:} Estamos en el último año de Ingeniería Informática en UADE y estamos haciendo la tesis. Desde la cátedra nos piden generar una validación con profesionales. Nuestra idea es desarrollar una herramienta de inteligencia artificial que, a partir de resonancias magnéticas, pueda detectar lesiones o generar mediciones y producir un preinforme para que después el profesional lo vea, lo valide y lo corrija si hace falta. Buscamos agilizar el proceso y automatizar tareas que pueden ser lentas o tediosas. La idea de estas preguntas es validar la propuesta.

\noindent\textbf{Dr.~Gonzalo Brito:} Bien. Entonces lo apuntan más al informe médico y no tanto a la adquisición de imágenes. En resonancia también se puede usar inteligencia artificial para adquirir secuencias, pero ustedes lo orientan más a pasar datos, analizarlos y apoyar el informe. Está bueno. La inteligencia artificial se empieza a ver y a escuchar mucho en nuestra especialidad porque trabajamos con computadoras, archivos y software. En nuestra empresa lo hemos usado en radiografía. Se testeó aplicarlo a placas simples y estaba bueno, pero tenía limitaciones: daba muchos falsos positivos y se empezó a dejar de usar. Habría que pulirlo. Igual, como preinforme o ayuda, puede servir. A veces se nos pueden escapar cosas y está bueno que un programa te marque una anormalidad o una imagen sospechosa para que después el médico la interprete y confirme si es patológica o si es un artefacto. En resonancia hay muchos artefactos, imágenes que parecen algo y probablemente no lo sean. La iniciativa es muy buena. Creo que a futuro se va a usar cada vez más y en congresos se habla mucho de esto.

\noindent\textbf{Dr.~Gonzalo Brito:} Al principio, hace algunos años, el miedo de los radiólogos era que la inteligencia artificial nos reemplazara. Pero ahora lo que se empieza a decir es que quizás la inteligencia artificial no va a reemplazar a los radiólogos, sino a los radiólogos que no se adapten a la inteligencia artificial. Si empieza a aparecer, ya no hay vuelta: está para quedarse. Probablemente los radiólogos que usen inteligencia artificial les saquen ventaja a los que no la usen.

\noindent\textbf{Francisco Fabrello:} Nosotros vamos por ese lado, que sea una herramienta para ayudar y no para reemplazar. Tampoco confiaríamos en que diga directamente «tenés tal cosa» hasta que lo revise una persona.

\noindent\textbf{Dr.~Gonzalo Brito:} Claro. Como algo previo lo veo útil. No creo que en este momento salga un informe completo sin que lo revise ninguna persona. Se ponen en juego muchas cosas. Pero sí creo que promete y que va a venir cada vez más. Habrá que pulirlo para que sea útil para nosotros y para el paciente. También habrá que estudiar si permite diagnosticar más ciertas patologías o detectar lesiones pequeñas que al ojo humano se le pueden pasar. Si la inteligencia artificial te marca una lesión de dos milímetros y después el médico decide si es relevante o no, puede ser muy útil.

\noindent\textbf{Francisco Fabrello:} Primero, para poner un poco el eje, ¿cómo sería tu trabajo con resonancia y qué lugar ocupa dentro de tu práctica?

\noindent\textbf{Dr.~Gonzalo Brito:} Ahora trabajo de lunes a jueves con resonancia magnética. Informo cuatro horas por día. A la mañana hago ecografías y, en la parte de imágenes, también informo tomografía. Hago bastante más tomografía, pero la resonancia me lleva un tiempo similar según la complejidad. Puedo hacer cinco o seis resonancias por hora, dependiendo del estudio. Hay estudios más largos, estudios oncológicos que requieren comparar y otros más sencillos que se ven más rápido.

\noindent\textbf{Francisco Fabrello:} ¿Cómo sería tu recorrido desde que revisás el estudio hasta que redactás el informe?

\noindent\textbf{Dr.~Gonzalo Brito:} Abro el estudio desde el sistema que tenemos, que se llama EGES. Ahí se organizan los pacientes, las órdenes y los estudios previos. Tener estudios previos es fundamental; siempre decimos que el mejor amigo del radiólogo son los estudios previos. Después abro el estudio, cargo las secuencias y empiezo a analizarlas. Luego grabo el informe con micrófono. El sistema permite grabar y eso va a un equipo de tipistas. Si digo que lo confirmen, el informe queda confirmado. Si no, queda sin confirmar y después entro manualmente, modifico lo que haga falta y cuando lo confirmo queda cargado en el sistema para que otros médicos lo vean en el portal.

\noindent\textbf{Francisco Fabrello:} ¿Lo hacés con varias vistas o con alguna vista particular? Por ejemplo, axial y sagital al mismo tiempo.

\noindent\textbf{Dr.~Gonzalo Brito:} Todas las reconstrucciones son fundamentales. Los cortes axiales suelen ser los que más usamos porque generalmente se planifica todo en cortes axiales y después se hacen reconstrucciones sagitales y coronales. En algunas situaciones específicas se hacen reconstrucciones oblicuas o siguiendo el eje de una estructura que quieras evaluar. Lo mejor es ver todos los planos porque te dan mucha información.

\noindent\textbf{Francisco Fabrello:} En columna, ¿qué hallazgo considerás más importante para que el informe sea clínicamente útil?

\noindent\textbf{Dr.~Gonzalo Brito:} Depende mucho del estudio, pero en columna lo que más busca el traumatólogo o el médico derivante es ver si están comprometidos los neuroforámenes y si las raíces nerviosas están comprometidas por algún disco. También puede haber patología oncológica, lesiones secundarias o hallazgos relevantes en columna. En general depende de qué estés buscando.

\noindent\textbf{Dr.~Gonzalo Brito:} En estudios donde importa lo oncológico o infeccioso, una secuencia que miro mucho es difusión y mapa de ADC. La difusión permite ver si hay restricción al movimiento de las moléculas de agua, lo que puede asociarse a patología neoplásica, infecciosa, abscesos o tumores. Es una secuencia que da mucha información y es de lo que menos me quiero dejar pasar.

\noindent\textbf{Francisco Fabrello:} ¿Qué error de una herramienta de este tipo sería clínicamente grave o inaceptable?

\noindent\textbf{Dr.~Gonzalo Brito:} Creo que lo más grave sería que se pase algo importante. Pienso en lo oncológico, por ejemplo. Hay criterios para ver si una patología está en remisión, progresando o igual. Si la herramienta no permite evaluar o comparar bien lesiones muy pequeñas, donde quizás un milímetro más es significativo, eso sería grave. Para mí lo más grave es que no detecte un hallazgo importante.

\noindent\textbf{Francisco Fabrello:} ¿Sería peor un falso negativo que un falso positivo?

\noindent\textbf{Dr.~Gonzalo Brito:} Sí. Un falso positivo, si lo marca, después lo desestimás. El problema es cuando algo no lo detecta. También nos pasó con un programa de radiografías: lo podían ver médicos derivantes en el portal, y a veces marcaba fractura donde no había fractura por una imagen superpuesta o un artefacto. Para el radiólogo no es tan grave porque lo revisa, pero si lo ve alguien que no es radiólogo y actúa en base a eso, puede ser grave. Lo más delicado es que lo usen personas que no son radiólogas y se queden con ese resultado sin revisión.

\noindent\textbf{Francisco Fabrello:} ¿Qué tendría que mostrar o permitir el sistema para que te resulte confiable, fácil de usar y no te agregue ruido en la rutina?

\noindent\textbf{Dr.~Gonzalo Brito:} Creo que sería útil que pueda medir las imágenes que se ven o sacar volúmenes. Eso podría ahorrar tiempo. Si marca una lesión, que la delimite y saque medidas en los tres ejes. Eso sería útil para agilizar. También debería ser algo sencillo.

\noindent\textbf{Francisco Fabrello:} Por ejemplo, que genere contornos superpuestos sobre la imagen original, que permita corrección manual o que muestre un indicador de confianza.

\noindent\textbf{Dr.~Gonzalo Brito:} Sí, eso podría servir. También algo que dé diagnósticos diferenciales podría estar bueno. Yo a veces saco una foto de una imagen, se la mando a ChatGPT y le pido diagnósticos diferenciales. Me tira opciones y después reviso bibliografía para ver cómo deberían verse. Muchas veces ayuda. Si el programa lo hiciera automáticamente, con una pestaña de diagnósticos potenciales, podría ser útil. No para quedarse con eso, sino como ayuda para que se escapen menos cosas o para recordar diagnósticos raros.

\noindent\textbf{Francisco Fabrello:} Para que el programa sea defendible desde el punto de vista clínico, ¿contra qué debería validarse? ¿Más del lado de datos o de profesionales?

\noindent\textbf{Dr.~Gonzalo Brito:} Creo que más por el lado de los profesionales. Que lo usen, que lo prueben los médicos. Si es una herramienta para médicos, la validación debería venir primero por ahí: que sea cómodo, práctico y que no cambie demasiado la rutina. A veces nos quieren imponer herramientas que terminan tardando más. Si me tarda un minuto más por informe, al final del día pesa. Tiene que ser práctico y no engorroso.

\noindent\textbf{Francisco Fabrello:} ¿Qué alcance te parece más defendible? ¿Qué agregarías, qué sacarías o qué sería una red flag?

\noindent\textbf{Dr.~Gonzalo Brito:} Me parece que está bueno que tenga mediciones, que detecte imágenes y que además pueda dar información. También podría incorporar una base teórica o enlazarse con sitios radiológicos. Nosotros usamos mucho Radiopaedia, que es como nuestra biblia: están muchos hallazgos y casos. Si pudiera conectarse o buscar ahí, estaría bueno.

\noindent\textbf{Dr.~Gonzalo Brito:} Como red flag, vuelvo a lo de los falsos negativos: que no detecte algo importante. También estaría bueno poder reportar errores. Cuando usábamos el programa de placas, tenía un QR para reportar errores. Eso era clave porque le servía a quienes desarrollaban el sistema para mejorarlo. Reportar en qué se equivoca el sistema es muy importante.

\noindent\textbf{Francisco Fabrello:} Eso sería todo. Como conclusión rápida, varias entrevistas van por el mismo lado: que no deje de detectar algo importante, que se pueda corregir, que no retrase el flujo y que ayude con los planos axiales.

\noindent\textbf{Dr.~Gonzalo Brito:} Cualquier otra duda me escriben. Si necesitan algo, no hay problema.

\noindent\textbf{Francisco Fabrello:} Por ahora estamos en una etapa teórica, pero seguramente cuando avancemos en el desarrollo notemos nuevas dudas y podamos volver a consultar.

\noindent\textbf{Dr.~Gonzalo Brito:} Sí, me escriben por mail o WhatsApp.

\noindent\textbf{Francisco Fabrello:} ¿Querés comentar algo más que no hayamos preguntado?

\noindent\textbf{Dr.~Gonzalo Brito:} No, creo que hablamos bastante. Como comentario, apoyo la iniciativa. Me parece algo buenísimo y creo que cada vez va a venir más. En un futuro no muy lejano seguramente vamos a estar trabajando con estos programas. Métanle, está bueno.

\noindent\textbf{Francisco Fabrello:} Muchas gracias.

\noindent\textbf{Dr.~Gonzalo Brito:} Éxitos. Cualquier cosa seguimos en contacto.
% SOURCE_NUMBERING_ANOMALY_PRESERVED: E.3 literal in DOCX Heading2
\phantomsection
\section*{E.3 Entrevista 3 (EN-03) - Dra. Claudia Cejas}
\addcontentsline{toc}{section}{E.3 Entrevista 3 (EN-03) - Dra. Claudia Cejas}
\label{anexo:entrevista-en03}

\noindent\textbf{Dra. Claudia Cejas:} Actualmente estoy como Presidenta de la Sociedad Argentina de Radiología y, dentro de las imágenes, me dedico a todo lo que es patología neuromuscular y de columna. Eso es lo que hago desde hace 15 años. Antes hacía todo lo que es región de cabeza y cuello. Bueno, ¿se quieren presentar?

\noindent\textbf{Francisco Fabrello:} Yo soy Francisco Fabrello y mi compañero es Enzo Asplanatti. Estamos cursando nuestro último año de Ingeniería Informática en UADE y estamos haciendo nuestra tesis sobre resonancias magnéticas, específicamente de columna lumbar. Empezamos este ciclo de entrevistas para validar un poco el proyecto y ver qué cambios hacen falta, porque es un área nueva para nosotros.

\noindent\textbf{Enzo Asplanatti:} La idea es desarrollar una herramienta acompañada con inteligencia artificial que permita detectar patrones a través de mediciones y generar un preinforme que después tenga validación profesional. Nosotros venimos a acompañar y automatizar parte del proceso, no a reemplazar el criterio médico.

\noindent\textbf{Dra. Claudia Cejas:} ¿Con respecto a la modalidad, están hablando de radiología, tomografía o resonancia?

\noindent\textbf{Enzo Asplanatti:} Por ahora vamos a trabajar con resonancias. Después, si se puede llevar a tomografía o radiografías, bienvenido, pero por ahora nos manejamos con resonancia.

\noindent\textbf{Dra. Claudia Cejas:} Perfecto. Resonancia es lo que se pide, así que está bueno. Los escucho: ¿qué necesitan de mí?

\noindent\textbf{Francisco Fabrello:} Tengo una serie de preguntas estructuradas. Para comenzar: ¿cómo es tu trabajo habitual con resonancias magnéticas y qué lugar ocupan dentro de tu práctica?

\noindent\textbf{Dra. Claudia Cejas:} Dentro de neuromuscular veo mucha resonancia de columna, así que es una de las áreas que evalúo en mi día a día. Estoy muchas horas en FLENI y parte de ese horario lo dedico a hacer informes de resonancia. Dentro de esas resonancias veo resonancia de columna.

\noindent\textbf{Francisco Fabrello:} Cuando abrís una resonancia, ¿cómo es tu recorrido habitual desde que empezás a revisar hasta que redactás el informe?

\noindent\textbf{Dra. Claudia Cejas:} Todos los radiólogos tenemos un algoritmo. En mi caso empiezo por el hueso. Evalúo la alineación de la columna, si está bien alineada o no. Luego me fijo cómo está el cuerpo vertebral, si tiene alguna patología, y después todo el aparato posterior. Luego paso a los discos y los veo uno por uno: la señal, la altura, si hay abombamientos, protrusiones, extrusiones y si generan compresión de raíces. Después paso a la médula. En lumbar se ve poco, pero en cervical y dorsal se ve mucho. Todo esto lo estudio en los planos que usamos: sagital y axial; si el paciente tiene escoliosis, también coronal. Después miro la parte posterior de la columna, músculos paravertebrales, ligamentos y estructuras alrededor, como riñones o aorta, porque a veces el paciente viene por una cosa y encontrás otra.

\noindent\textbf{Francisco Fabrello:} ¿Eso lo hacés al mismo tiempo con dos pantallas o separando la imagen en un monitor?

\noindent\textbf{Dra. Claudia Cejas:} Uso dos pantallas. Divido cada pantalla en cuatro y voy cargando las imágenes. Como tenemos un sistema PACS, cuando estoy en el sagital puedo ver en qué nivel estoy. Entonces voy jugando para ver si la afectación está dentro de los neuroforámenes, en la parte superior, media o inferior.

\noindent\textbf{Francisco Fabrello:} ¿Qué hallazgos considerás más importantes para el informe que haría este sistema? Por ejemplo, degeneración discal, hernias, protrusiones, estenosis, canal espinal o forámenes.

\noindent\textbf{Dra. Claudia Cejas:} Te puedo hacer un listado. Pensando en mi algoritmo, primero la alineación de la columna, porque muchas veces en un informe previo no se informó que estaba desalineada. Eso es importante porque habla de un problema funcional. Después vería la presencia de protrusiones, sus características y si son protrusiones o extrusiones. Luego iría al grado de estenosis, ya sea del foramen, del espacio lateral o del canal central. Tener esas medidas estaría bueno porque nosotros muchas veces lo graduamos a ojo: leve, moderado o severo. Necesitaría un sistema para poder protocolizarlo.

\noindent\textbf{Dra. Claudia Cejas:} También detectaría lesiones en la médula: lesiones desmielinizantes o tumorales. Si vamos a columna degenerativa, me enfocaría en discos; pero si pensamos en un centro neurológico, una lesión en la médula no debería omitirse. También podría estar bueno detectar lesiones óseas, no solo en cuerpos vertebrales sino en pedículos y láminas. Las láminas se pueden romper y generar inestabilidad, y muchas veces eso nos cuesta verlo en resonancia.

\noindent\textbf{Francisco Fabrello:} Si el sistema trabajara solo con cortes sagitales y no con cortes axiales, ¿se podría evaluar razonablemente?

\noindent\textbf{Dra. Claudia Cejas:} Depende del tipo de secuencia. Si tenés secuencias volumétricas 3D, podés programar una secuencia en el plano sagital y reconstruirla en coronal y axial con la misma resolución. Eso es válido. Ahora, hacer solamente un corte sagital 2D, no. Te quedás corto. Es raro que haya lugares que hagan un solo plano.

\noindent\textbf{Francisco Fabrello:} Con las entrevistas nos quedó más claro que el axial es necesario para un buen análisis.

\noindent\textbf{Dra. Claudia Cejas:} Sí. Si tenés secuencias 3D podés reconstruir, pero si no, el axial es importante.

\noindent\textbf{Francisco Fabrello:} ¿Qué error de una herramienta de este tipo sería clínicamente grave o inaceptable?

\noindent\textbf{Dra. Claudia Cejas:} Si ustedes incluyen médula, lo más grave sería no ver una lesión en la médula. Si tenés una lesión en la región cervical, una persona puede quedar parapléjica o sin mover brazos y piernas. Creo que eso es lo que no debería perderse. Todo lo demás es más perdonable.

\noindent\textbf{Francisco Fabrello:} ¿Qué podría mostrar o permitir el sistema para que te resulte confiable y mejore la rutina?

\noindent\textbf{Dra. Claudia Cejas:} Que me marque con algún color. Por ejemplo, que me delinee el disco. También me imagino que me delinee el borde posterior de la columna en sagital para ver si hay desviación. Eso me encantaría. Si me delinea el disco, puedo ver rápidamente si hay solo un abombamiento o si el disco coincide con el cuerpo vertebral, que sería normal. Si está por fuera o tiene una deformidad, me llama la atención para ver una protrusión. También me serviría que me marque la médula o alguna lesión.

\noindent\textbf{Francisco Fabrello:} La idea sería tener la resonancia y aplicar una máscara arriba, pintando las partes. También pensamos en que marque un nivel de confianza del resultado, por ejemplo, una medida con cierto porcentaje de confianza.

\noindent\textbf{Dra. Claudia Cejas:} Sí, eso estaría bueno. ¿Y la idea de ustedes es hacer un reporte agregado a eso?

\noindent\textbf{Francisco Fabrello:} Sí. Una idea era que oriente al médico con posibles casos o diagnósticos diferenciales, pero sin decir automáticamente «tiene esto». La idea es ayudar a plantear opciones para que el médico llegue a la conclusión.

\noindent\textbf{Dra. Claudia Cejas:} Yo creo que lo ideal, para acelerar el trabajo hoy, es que la herramienta esté destinada a la detección, no a decir qué es. Para llegar a decir qué es nos basamos mucho en la clínica. Una lesión sola puede no coincidir con la clínica. En un hospital, además, uno lee la historia clínica y eso cambia el diagnóstico. Hoy no esperaría que una inteligencia artificial inicial me diga el diagnóstico. Esperaría que me detecte lesiones, que me muestre lo que yo sé ver pero que podría omitir por carga laboral o burnout. Me parece que ese debería ser el enfoque. El diagnóstico con clínica, laboratorio y todo lo demás puede venir en una segunda etapa.

\noindent\textbf{Dra. Claudia Cejas:} Nosotros trabajamos con inteligencia artificial en neuro. Validamos estudios de una empresa que detecta atrofia cerebral y lesiones desmielinizantes cerebrales. Nos hace un reporte estructurado y después nosotros vemos si ese reporte queda. En atrofia es muy útil porque antes lo hacíamos a ojo y había mucha discordancia. El sistema mide atrofia global y por áreas específicas. Eso acelera muchísimo.

\noindent\textbf{Francisco Fabrello:} Para que el producto sea defendible desde el punto de vista clínico, ¿contra qué considerás que deberíamos validarlo?

\noindent\textbf{Dra. Claudia Cejas:} Con expertos. Si es un producto que no está desarrollado, lo validaría con un grupo de expertos.

\noindent\textbf{Francisco Fabrello:} Nos habían mencionado Radiopaedia como referencia.

\noindent\textbf{Dra. Claudia Cejas:} Radiopaedia está buena. Habla de la patología y al costado pone casos clínicos. La usan mucho los residentes para buscar diagnósticos diferenciales. Está muy bien hecha.

\noindent\textbf{Francisco Fabrello:} ¿Qué alcance te parece más defendible? ¿Qué sacarías para no prometer de más?

\noindent\textbf{Dra. Claudia Cejas:} Lo que pretendería es que puedan dar toda la información que les mencioné, pero no sé si es viable. Tendrían que limitarse a algo. Si es un centro neurológico como el nuestro, algo que no debería omitirse es una lesión en la médula. Ahora, en un centro de diagnóstico, donde la patología degenerativa es lo más común en columna, me dedicaría a los discos.

\noindent\textbf{Dra. Claudia Cejas:} Otra opción sería medir el canal. Si el canal está estrecho, quien lo ve va a buscar por qué está estrecho. Puede ser porque está desalineado, porque tiene estenosis o porque tiene un tumor. Para simplificarlo, podrían medir el canal lumbar. Está estandarizado en bibliografía y saben cuánto tiene que medir en un adulto. Si lo hacen, háganlo en adultos, porque pediatría es otro mundo. Acótenlo a región lumbar y a un tema específico.

\noindent\textbf{Francisco Fabrello:} ¿Querés comentarnos algo más?

\noindent\textbf{Dra. Claudia Cejas:} Me parece hermoso el proyecto que van a armar. Está buenísimo. ¿Qué herramientas van a usar? ¿Python, R? ¿Cómo van a desarrollarlo?

\noindent\textbf{Francisco Fabrello:} Para la parte de IA pensamos en Python. Para la interfaz estamos viendo algo más web, para que no tenga que correr localmente ni requiera un equipo potente. La idea es que el usuario cargue el estudio y se procese en la nube.

\noindent\textbf{Dra. Claudia Cejas:} Ah, ok.

\noindent\textbf{Francisco Fabrello:} También pensamos en la posibilidad de cargar una resonancia anterior y la actual para que haga una comparativa: qué cambió, qué no y qué se mantuvo.

\noindent\textbf{Dra. Claudia Cejas:} Eso es espectacular. No sé si es sencillo de hacer, pero es espectacular, porque nos pasa mucho que piden controles y tenemos que ir viendo hacia atrás.

\noindent\textbf{Francisco Fabrello:} Vamos a ver si es técnicamente viable.

\noindent\textbf{Dra. Claudia Cejas:} Nosotros tenemos un sistema que entra al PACS, lee el estudio, lo analiza y nos lo devuelve. También conozco otros sistemas en la nube donde cargás el estudio, lo analiza y lo devuelve. También está bueno.

\noindent\textbf{Francisco Fabrello:} Estamos yendo más hacia la nube porque correr IA localmente requiere muchos recursos. Estamos haciendo entrenamientos preliminares y creemos que puede funcionar.

\noindent\textbf{Dra. Claudia Cejas:} Bueno, mucha suerte.

\noindent\textbf{Francisco Fabrello:} ¿Te podríamos contactar en una futura entrevista para mostrarte cómo avanzamos y recibir feedback?

\noindent\textbf{Dra. Claudia Cejas:} Sí, claro que sí. No hay problema.

\noindent\textbf{Francisco Fabrello:} Muchas gracias.

\noindent\textbf{Dra. Claudia Cejas:} Gracias a ustedes. Que tengan buen día.
% MIGRATION_PENDING: E.4 / EN-04
% MIGRATION_PENDING: E.5 / EN-05
